% Propuesta de inserción después de 07b_geometria_elipse.tex.
% RESULTADO_RECUPERADO: c52_06_01_incertidumbre.tex, líneas 10--91.
% Dependencias: c52_04_01_ley_areal.tex, líneas 137--198;
% c52_06_02_intercambio_constitutivo.tex, líneas 5--29.
% Explicitación local: dominio operatorio, prueba de covarianza,
% transporte recíproco y distinción dimensional posición--momento.
% Este fragmento no introduce una constante ni un generador adicional.

\subsection{Covariance canonique et réciprocité de l'ellipse d'action}
\label{sec:ii-area-covarianza}

La construction précédente détermine deux invariants complémentaires :
le produit des demi-axes fixe l’aire d’action et leur quotient conserve
l’anisotropie de la paire angulaire. Précisons maintenant la réalisation canonique
dans laquelle ces demi-axes sont deux écarts-types. La structure de
départ est l’ellipse déjà obtenue à partir des représentations angulaires et
d’action d’APP--TRIT--TPK ; la réalisation opératorielle s’établit après
cette construction.

Dans ce paragraphe, on écrit \(\eta=\eta_{\rm el}
=\operatorname{artanh}(C^*/A)\), avec \(|C^*/A|<1\) et
\(\hbar>0\). Le quotient angulaire utilise la même unité pour ses deux
composantes. Dans les coordonnées équilibrées \((Q,P)\), toutes deux de
dimension racine carrée d'action, définissons
\begin{equation}
 V_\eta:=\frac14
 \begin{pmatrix}a_S^2&0\\0&b_S^2\end{pmatrix}
 =\frac{\hbar}{2}
 \begin{pmatrix}e^\eta&0\\0&e^{-\eta}\end{pmatrix}.
 \label{eq:ii-covarianza-geometrica}
\end{equation}
Cette matrice positive possède des entrées de dimension action. Son interprétation
comme covariance quantique est spécifiée par la réalisation suivante.

\begin{proposition}[Réalisation gaussienne de la covariance d'action]
\label{prop:ii-covarianza-gaussiana}
Considérons la représentation canonique dans \(L^2(\mathbb R,dQ)\), avec
\(\widehat Qf(Q)=Qf(Q)\) et
\(\widehat Pf(Q)=-i\hbar f'(Q)\), initialement sur l'espace de
Schwartz \(\mathcal S(\mathbb R)\). L'état normalisé
\begin{equation}
 \psi_\eta(Q)=\bigl(\pi\hbar e^\eta\bigr)^{-1/4}
 \exp\!\left(-\frac{Q^2}{2\hbar e^\eta}\right)
 \label{eq:ii-gaussiana-accion}
\end{equation}
a des moyennes nulles et pour matrice de covariance \(V_\eta\). En particulier,
\begin{equation}
 (\Delta Q)^2=\frac{\hbar e^\eta}{2},\qquad
 (\Delta P)^2=\frac{\hbar e^{-\eta}}{2},\qquad
 \operatorname{Cov}(Q,P)=0,\qquad
 \det V_\eta=\frac{\hbar^2}{4}.
 \label{eq:ii-varianzas-accion}
\end{equation}
Cette réalisation sature l'inégalité de Robertson--Schrödinger.
\end{proposition}

\begin{proof}
Les opérateurs précédents sont essentiellement autoadjoints sur
\(\mathcal S(\mathbb R)\), qui constitue un cœur commun invariant ;
sur cet espace, on a
\([\widehat Q,\widehat P]=i\hbar I\). Pour
\(\widehat Y=(\widehat Q,\widehat P)\), la covariance symétrique est
\[
 V_{jk}=\frac12\left\langle
 (\widehat Y_j-\langle\widehat Y_j\rangle)
 (\widehat Y_k-\langle\widehat Y_k\rangle)
 +(\widehat Y_k-\langle\widehat Y_k\rangle)
 (\widehat Y_j-\langle\widehat Y_j\rangle)
 \right\rangle.
\]
Avec \(s=\hbar e^\eta>0\), soit
\(I_s=\int_{\mathbb R}e^{-Q^2/s}\,dQ\). Le produit de deux
intégrales identiques, évalué en coordonnées polaires, donne
\(I_s^2=2\pi\int_0^\infty r e^{-r^2/s}\,dr=\pi s\).
En outre, l'intégrale de
\(\frac{d}{dQ}(Qe^{-Q^2/s})\) est nulle, par décroissance aux deux
extrémités ; d'où
\(\int_{\mathbb R}Q^2e^{-Q^2/s}\,dQ=sI_s/2\).
Ces identités et la parité démontrent la normalisation, la moyenne nulle et
\(\langle\widehat Q^2\rangle=s/2\). De plus,
\(\widehat P\psi_\eta=i e^{-\eta}Q\psi_\eta\).
Ainsi, \(\langle\widehat P\rangle=0\),
\(\|\widehat P\psi_\eta\|^2=\hbar e^{-\eta}/2\) et
\(\operatorname{Re}\langle\widehat Q\psi_\eta,
\widehat P\psi_\eta\rangle=0\).

Pour tout état normalisé de ce domaine, Cauchy--Schwarz appliquée
à \((\widehat Q-\langle\widehat Q\rangle)\psi\) et
\((\widehat P-\langle\widehat P\rangle)\psi\) donne
\[
 (\Delta Q)^2(\Delta P)^2
 \geq \operatorname{Cov}(Q,P)^2
       +\frac14\left|\langle[\widehat Q,\widehat P]\rangle\right|^2
 =\operatorname{Cov}(Q,P)^2+\frac{\hbar^2}{4}.
\]
Les variances calculées atteignent l'égalité. De manière équivalente,
l'opérateur
\[
 a_\eta=\frac{e^{-\eta/2}\widehat Q
                  +i e^{\eta/2}\widehat P}{\sqrt{2\hbar}}
\]
satisfait \([a_\eta,a_\eta^\dagger]=I\) et
\(a_\eta\psi_\eta=0\), par substitution directe.
\end{proof}

\begin{proposition}[Aire, dispersion et transport réciproque]
\label{prop:ii-area-covarianza-reciprocidad}
L'ellipse d'action est exactement le sous-niveau de covariance
\begin{equation}
 \mathcal E_\eta
 =\left\{y=(Q,P)^{\mathsf T}:y^{\mathsf T}V_\eta^{-1}y\leq4\right\}
 =\left\{\frac{Q^2}{a_S^2}+\frac{P^2}{b_S^2}\leq1\right\}.
 \label{eq:ii-elipse-covarianza}
\end{equation}
Son aire est \(h=2\pi\hbar\), ses demi-axes sont
\(a_S=2\Delta Q\), \(b_S=2\Delta P\) et
\begin{equation}
 \frac{\Delta P}{\Delta Q}=e^{-\eta}
 =\frac{b_S}{a_S}=R_{\rm el}^2.
 \label{eq:ii-covarianza-forma}
\end{equation}
La rotation canonique \(J_2\) de la section~\ref{sec:ellipse} transporte
\(V_\eta\) vers \(V_{-\eta}\) et conserve la forme symplectique.
\end{proposition}

\begin{proof}
L'inversion de la matrice diagonale
\eqref{eq:ii-covarianza-geometrica} donne
\eqref{eq:ii-elipse-covarianza}. Ses demi-axes sont les racines positives de
quatre fois ses variances, et
\[
 \operatorname{Area}(\mathcal E_\eta)
 =4\pi\sqrt{\det V_\eta}=2\pi\hbar=h,
\]
en accord avec la loi aréale~\eqref{eq:ii-elipse-area-fase}.
Le rapport des racines positives démontre
\eqref{eq:ii-covarianza-forma}.
Si \(M_\eta=\operatorname{diag}(e^{\eta/2},e^{-\eta/2})\), alors
\[
 V_\eta=M_\eta V_0M_\eta^{\mathsf T},\qquad
 M_\eta^{\mathsf T}J_2M_\eta=J_2,\qquad
 J_2V_\eta J_2^{\mathsf T}=V_{-\eta}.
\]
Ces identités s'obtiennent par multiplication. En termes opératoires,
\((\widehat Q',\widehat P')=(-\widehat P,\widehat Q)\) conserve
\([\widehat Q',\widehat P']=i\hbar I\). La permutation
\(S_2:(Q,P)\mapsto(P,Q)\) échange également les variances, mais inverse
le signe du commutateur et de la forme symplectique. Le transport de
covariance respecte ainsi la distinction entre réciprocité canonique et inversion
d'orientation déjà établie à la section~\ref{sec:ellipse}.
\end{proof}

\paragraph{Position, quantité de mouvement et transduction dimensionnelle.}
Soit \(\lambda_{xp}>0\) un facteur de la coordinatisation dimensionnelle, de dimension
quantité de mouvement par unité de longueur. La transformation
\[
 Q=\sqrt{\lambda_{xp}}\,x,\qquad
 P=\frac{p_x}{\sqrt{\lambda_{xp}}}
\]
réalise le couple équilibré dans une position \(x\) et sa quantité de mouvement conjuguée
\(p_x\) ; elle conserve \(dQ\wedge dP=dx\wedge dp_x\) et
\([\widehat x,\widehat p_x]=i\hbar I\). Par transformation de la
covariance,
\begin{equation}
 V_\eta^{(x,p_x)}=\frac{\hbar}{2}
 \begin{pmatrix}e^\eta/\lambda_{xp}&0\\
                 0&\lambda_{xp}e^{-\eta}\end{pmatrix},\qquad
 \Delta x\,\Delta p_x=\frac{\hbar}{2},\qquad
 \frac{\Delta p_x}{\Delta x}=\lambda_{xp}e^{-\eta}.
 \label{eq:ii-covarianza-dimensiones}
\end{equation}
Le quotient de dispersion acquiert ici la dimension quantité de mouvement par longueur.
En revanche, \eqref{eq:ii-covarianza-forma} appartient à la coordinatisation équilibrée.
Pour la réalisation inductive--capacitive déjà définie,
\eqref{eq:ii-elipse-impedancia} apporte l’identité adimensionnelle
\(Z_\eta/Z_*=e^{-\eta}=\Delta P/\Delta Q\).
Chaque lecture utilise sa transduction déclarée et conserve le même
déterminant d’action.

\paragraph{Énergie du contour et énergie moyenne de l'état.}
La réalisation inductive--capacitive possède le hamiltonien
\[
 H_\eta(Q,P)=\frac{\omega}{2}
 \left(e^{-\eta}Q^2+e^\eta P^2\right),\qquad \omega>0.
\]
Sur \(\partial\mathcal E_\eta\), sa valeur est
\(H_\eta=\hbar\omega\), d'après les équations des demi-axes. Dans la
réalisation quantique de la proposition~\ref{prop:ii-covarianza-gaussiana},
\begin{equation}
 \langle\widehat H_\eta\rangle_{\psi_\eta}
 =\frac{\omega}{2}
 \left(e^{-\eta}\frac{\hbar e^\eta}{2}
       +e^\eta\frac{\hbar e^{-\eta}}{2}\right)
 =\frac{\hbar\omega}{2}.
 \label{eq:ii-energia-contorno-covarianza}
\end{equation}
Le contour de deux écarts-types possède donc le double de
l'énergie moyenne de cet état gaussien. La loi aréale fixe la normalisation
d'action ; le couple angulaire fixe sa distribution entre les deux coordonnées ;
la représentation canonique et l'état explicite réalisent cette distribution
comme covariance. La saturation de l'incertitude est ainsi établie dans
la réalisation spécifiée par ses opérateurs et son état.
